\section{Questao 3 - Busca Local (8-Rainhas)}

\subsection{Formulacao em espaco de estados}
\[
P = \langle Z, z_0, \text{Goal}, \text{Succ}, c \rangle
\]
\begin{itemize}[leftmargin=1.5em]
  \item a) Estado formal: $z=(s_1,\dots,s_8)$, com $s_i \in \{1,\dots,8\}$ indicando a linha da rainha na coluna $i$. Assim, $Z=\{1,\dots,8\}^8$.
  \item b) Estado inicial: $z_0=(1,1,1,1,1,1,1,1)$.
  \item c) Teste de objetivo: $\text{Goal}(z)$ verdadeiro sse $h(z)=0$, onde $h(z)$ e o numero de pares de rainhas em conflito (mesma linha ou diagonal).
  \item d) Funcao sucessora: $\text{Succ}(z)$ gera todos os estados obtidos movendo exatamente uma rainha para outra linha na mesma coluna.
  \item e) Funcao de custo: $c(z,z')=1$ por movimento local (vizinho); em busca local, minimiza-se principalmente $h(z)$.
  \item f) Representacao em Python: tupla de 8 inteiros (ex.: \texttt{(1,1,1,1,1,1,1,1)}).
\end{itemize}

\textbf{Interpretacao da superficie de busca (local):}
\begin{itemize}[leftmargin=1.5em]
  \item A paisagem de $h(s)$ e discreta e multimodal, com muitos planaltos e minimos locais.
  \item Movimentos de uma unica rainha por coluna podem reduzir conflitos localmente, mas tambem podem levar a regioes sem melhoria estrita.
\end{itemize}

\subsection{Hill-Climbing (execucao manual)}
\begingroup
\small
\begin{longtable}{c p{3.8cm} c p{7.2cm}}
\toprule
Iteracao & Estado atual & $h(s)$ & Estado escolhido e justificativa \\
\midrule
0 & [1,1,1,1,1,1,1,1] & 28 & [1,1,1,1,1,1,8,1] (h=21): melhor vizinho \\
1 & [1,1,1,1,1,1,8,1] & 21 & [1,1,1,1,1,1,8,2] (h=15): melhor vizinho \\
2 & [1,1,1,1,1,1,8,2] & 15 & [1,1,1,1,7,1,8,2] (h=10): melhor vizinho \\
3 & [1,1,1,1,7,1,8,2] & 10 & [1,1,6,1,7,1,8,2] (h=6): melhor vizinho \\
4 & [1,1,6,1,7,1,8,2] & 6 & [5,1,6,1,7,1,8,2] (h=3): melhor vizinho \\
5 & [5,1,6,1,7,1,8,2] & 3 & [5,1,6,4,7,1,8,2] (h=1): melhor vizinho \\
6 & [5,1,6,4,7,1,8,2] & 1 & parada: nao existe vizinho com $h$ menor \\
\bottomrule
\end{longtable}
\endgroup

\textbf{Resultado do HC manual:}
\begin{itemize}[leftmargin=1.5em]
  \item Estado final: [5,1,6,4,7,1,8,2].
  \item $h$ final: 1.
  \item Nao encontrou solucao valida ($h \neq 0$).
  \item Numero total de iteracoes executadas: 7 (iteracoes 0 a 6), com 6 melhorias e 1 parada.
  \item Tipo de parada: \textbf{plato} (existem vizinhos com mesmo valor de $h$, mas nenhum melhor).
\end{itemize}

\subsection{Top-5 vizinhos por iteracao (HC)}
\begingroup
\scriptsize
\begin{longtable}{c p{4.7cm} c}
\toprule
Iteracao & Vizinho & $h(s)$ \\
\midrule
0 & [1,1,1,1,1,1,8,1] & 21 \\
0 & [1,1,1,1,1,7,1,1] & 21 \\
0 & [1,1,1,1,1,8,1,1] & 21 \\
0 & [1,1,1,1,6,1,1,1] & 21 \\
0 & [1,1,1,1,7,1,1,1] & 21 \\
1 & [1,1,1,1,1,1,8,2] & 15 \\
1 & [1,1,1,1,7,1,8,1] & 15 \\
1 & [1,1,1,6,1,1,8,1] & 15 \\
1 & [1,1,1,7,1,1,8,1] & 15 \\
1 & [1,1,5,1,1,1,8,1] & 15 \\
2 & [1,1,1,1,7,1,8,2] & 10 \\
2 & [1,1,1,7,1,1,8,2] & 10 \\
2 & [1,1,5,1,1,1,8,2] & 10 \\
2 & [1,1,6,1,1,1,8,2] & 10 \\
2 & [1,6,1,1,1,1,8,2] & 10 \\
3 & [1,1,6,1,7,1,8,2] & 6 \\
3 & [1,6,1,1,7,1,8,2] & 6 \\
3 & [5,1,1,1,7,1,8,2] & 6 \\
3 & [1,1,1,1,7,2,8,2] & 7 \\
3 & [1,1,1,1,7,3,8,2] & 7 \\
4 & [5,1,6,1,7,1,8,2] & 3 \\
4 & [1,1,6,1,7,2,8,2] & 4 \\
4 & [1,1,6,1,7,4,8,2] & 4 \\
4 & [1,1,6,1,7,5,8,2] & 4 \\
4 & [1,1,6,2,7,1,8,2] & 4 \\
5 & [5,1,6,4,7,1,8,2] & 1 \\
5 & [5,1,6,1,7,2,8,2] & 2 \\
5 & [5,1,6,1,7,4,8,2] & 2 \\
5 & [5,2,6,1,7,1,8,2] & 2 \\
5 & [5,1,6,1,3,1,8,2] & 3 \\
6 & [5,1,6,4,7,3,8,2] & 1 \\
6 & [5,3,6,4,7,1,8,2] & 1 \\
6 & [3,1,6,4,7,1,8,2] & 2 \\
6 & [5,1,6,4,3,1,8,2] & 2 \\
6 & [5,1,6,4,7,1,3,2] & 2 \\
\bottomrule
\end{longtable}
\endgroup

\subsection{Todos os vizinhos avaliados por iteracao (HC)}
\input{results/q3_hc_todos_vizinhos.tex}

\subsection{Analise de armadilha de busca}
\begin{itemize}[leftmargin=1.5em]
  \item Tipo observado: \textbf{plato} (na iteracao final, melhor vizinho possui o mesmo valor de $h$).
  \item Por que ocorreu: a vizinhanca de um movimento por coluna cria regioes planas no espaco de busca.
  \item Impacto: o HC puro para sem atingir $h=0$, mesmo estando proximo da solucao.
  \item Como evitar: random restart, movimentos laterais limitados ou simulated annealing.
\end{itemize}

\subsection{Random Restart Hill-Climbing (20 execucoes)}
\begingroup
\scriptsize
\begin{longtable}{c p{3.5cm} c c c}
\toprule
Execucao & Estado inicial & Passos & $h(s)$ final & Solucao? \\
\midrule
1 & [2,6,2,4,7,8,8,4] & 4 & 1 & nao \\
2 & [1,2,2,5,2,8,1,8] & 3 & 1 & nao \\
3 & [6,4,7,5,6,6,7,2] & 4 & 1 & nao \\
4 & [6,2,5,5,8,3,5,1] & 1 & 3 & nao \\
5 & [6,6,8,7,2,7,8,2] & 5 & 0 & sim \\
6 & [7,8,7,5,7,8,1,4] & 3 & 1 & nao \\
7 & [3,1,2,7,8,3,8,4] & 3 & 0 & sim \\
8 & [2,8,8,5,7,5,8,7] & 4 & 1 & nao \\
9 & [7,4,5,5,1,1,3,8] & 2 & 1 & nao \\
10 & [8,6,4,5,8,8,8,7] & 3 & 1 & nao \\
11 & [1,5,6,8,2,8,5,8] & 3 & 1 & nao \\
12 & [3,7,8,2,5,8,6,2] & 2 & 2 & nao \\
13 & [4,7,3,8,1,2,1,8] & 2 & 1 & nao \\
14 & [2,4,4,2,1,5,6,3] & 4 & 0 & sim \\
15 & [1,6,8,3,4,8,3,5] & 1 & 1 & nao \\
16 & [2,7,3,3,5,7,7,3] & 4 & 2 & nao \\
17 & [4,7,1,2,7,2,1,6] & 3 & 1 & nao \\
18 & [6,1,1,8,6,5,8,6] & 5 & 0 & sim \\
19 & [2,8,8,3,3,4,7,6] & 3 & 0 & sim \\
20 & [8,4,2,7,1,4,3,1] & 3 & 1 & nao \\
\bottomrule
\end{longtable}
\endgroup

\textbf{Resumo RR-HC:}
\begin{itemize}[leftmargin=1.5em]
  \item Taxa de sucesso: $5/20 = 25.0\%$.
  \item Media de $h$ final: 0.95.
  \item Media de passos por execucao: 3.10.
\end{itemize}

\subsection{Simulated Annealing}
\[
P(\text{aceitar pior}) = e^{-\Delta E / T}
\]
\begin{itemize}[leftmargin=1.5em]
  \item Temperatura inicial: $T_0 = 25.0$.
  \item Resfriamento: $T \leftarrow \alpha T$ com $\alpha = 0.97$.
  \item Parada: $T < 0.001$ ou limite de 3000 passos.
  \item Execucao detalhada: inicio [1,6,8,1,5,2,8,8], melhor estado [3,6,2,7,1,4,8,5], \textbf{best $h=0$} (sucesso), 311 passos.
\end{itemize}

\textbf{Exemplos de movimentos piores aceitos (execucao detalhada):}
\begingroup
\scriptsize
\begin{table}[H]
\centering
\begin{tabular}{c c c c c c}
\toprule
Step & $h_{atual}\rightarrow h_{novo}$ & $\Delta$ & $T$ & $e^{-\Delta/T}$ & sorteio \\
\midrule
5  & $5\rightarrow8$  & +3 & 21.4684 & 0.8696 & 0.3689 \\
6  & $8\rightarrow12$ & +4 & 20.8243 & 0.8252 & 0.4443 \\
13 & $6\rightarrow7$  & +1 & 16.8257 & 0.9423 & 0.8105 \\
14 & $7\rightarrow8$  & +1 & 16.3209 & 0.9406 & 0.4154 \\
19 & $4\rightarrow5$  & +1 & 14.0153 & 0.9311 & 0.6297 \\
22 & $4\rightarrow5$  & +1 & 12.7914 & 0.9248 & 0.3176 \\
23 & $5\rightarrow7$  & +2 & 12.4077 & 0.8511 & 0.1797 \\
24 & $7\rightarrow9$  & +2 & 12.0354 & 0.8469 & 0.3187 \\
\bottomrule
\end{tabular}
\end{table}
\endgroup

\textbf{Lote SA (20 execucoes):}
\begingroup
\scriptsize
\begin{longtable}{c p{3.5cm} c c c}
\toprule
Execucao & Estado inicial & Passos & best $h$ & Solucao? \\
\midrule
1 & [5,1,6,3,4,6,1,2] & 333 & 1 & nao \\
2 & [7,8,4,5,7,1,7,4] & 333 & 1 & nao \\
3 & [6,3,1,7,7,8,1,7] & 291 & 0 & sim \\
4 & [2,1,3,8,7,5,7,2] & 173 & 0 & sim \\
5 & [4,5,3,8,7,3,7,1] & 333 & 1 & nao \\
6 & [8,8,7,2,2,6,6,1] & 333 & 1 & nao \\
7 & [2,1,8,3,4,5,1,2] & 333 & 1 & nao \\
8 & [1,2,5,2,5,5,6,2] & 333 & 1 & nao \\
9 & [2,8,8,3,3,6,6,6] & 319 & 0 & sim \\
10 & [1,3,2,7,1,7,5,6] & 333 & 1 & nao \\
11 & [2,6,5,2,3,5,3,2] & 333 & 1 & nao \\
12 & [6,4,8,1,5,7,6,2] & 333 & 1 & nao \\
13 & [2,2,7,1,2,6,2,5] & 333 & 1 & nao \\
14 & [4,5,2,7,4,7,6,3] & 110 & 0 & sim \\
15 & [7,8,7,3,4,5,4,6] & 333 & 1 & nao \\
16 & [4,1,8,8,4,7,2,8] & 333 & 1 & nao \\
17 & [5,6,5,5,1,4,4,4] & 333 & 1 & nao \\
18 & [8,7,2,6,4,3,4,6] & 333 & 1 & nao \\
19 & [2,8,5,4,7,3,8,8] & 333 & 1 & nao \\
20 & [2,2,6,7,3,1,8,6] & 333 & 1 & nao \\
\bottomrule
\end{longtable}
\endgroup

\textbf{Resumo SA:}
\begin{itemize}[leftmargin=1.5em]
  \item Taxa de sucesso: $4/20 = 20.0\%$.
  \item Media de best $h$: 0.80.
  \item Media de passos por execucao: 311.05.
  \item Comparacao direta com HC simples: o HC manual parou em $h=1$, enquanto o SA alcancou $h=0$ na execucao detalhada e encontrou 4 solucoes validas em 20 execucoes.
\end{itemize}

\subsection{Comparacao final (Q3)}
\begingroup
\footnotesize
\begin{table}[H]
\centering
\setlength{\tabcolsep}{3pt}
\renewcommand{\arraystretch}{1.15}
\begin{tabular}{>{\raggedright\arraybackslash}p{3.8cm} >{\raggedright\arraybackslash}p{3.2cm} >{\raggedright\arraybackslash}p{3.2cm} >{\raggedright\arraybackslash}p{3.2cm}}
\toprule
Criterio & Hill-Climbing & RR Hill-Climbing & Simulated Annealing \\
\midrule
Qualidade da solucao & Falhou no caso manual ($h=1$) & 25\% de sucesso (5/20) & 20\% de sucesso (4/20), com melhor best $h$ medio (0.80) \\
Velocidade de convergencia & Muito rapido (6 melhorias) & Rapido por tentativa (3.10 passos em media) & Mais lento (311.05 passos em media) \\
Sensibilidade ao estado inicial & Alta & Media & Media/baixa \\
Escapa de maximos locais/platos & Nao & Parcialmente (por reinicio) & Sim (probabilistico) \\
Custo computacional & Baixo & Medio & Alto \\
Estabilidade dos resultados & Deterministico para um estado inicial fixo & Variavel entre execucoes & Variavel e sensivel aos parametros de temperatura \\
\bottomrule
\end{tabular}
\caption{Comparacao dos metodos de busca local para 8-rainhas}
\end{table}
\endgroup
